% Tracce uscite agli Esami di Stato (Sezione A, Settore Informazione) sugli argomenti del capitolo.
% Incluso da main.tex subito dopo il riassunto.

\section*{Tracce Uscite} % Usiamo * per non numerare la sezione nell'indice
\addcontentsline{toc}{section}{Tracce Uscite: Sistemi Distribuiti}

\subsection*{Prima sessione 2024 -- Informatica \#1}
\textbf{Traccia}: consistenza del dato in un sistema distribuito; soluzioni tipiche di database SQL e NoSQL con esempi e casi in cui un'azienda preferirebbe l'uno o l'altro; teorema CAP e il suo effetto sulla scelta.

\textbf{Dove trovare la risposta}: sezioni \emph{Il Teorema CAP} e \emph{Database SQL e NoSQL}; scaletta a fine sezione.

\subsection*{Prima sessione 2024 -- Informatica \#2}
\textbf{Traccia}: principali caratteristiche di un sistema distribuito e tecniche che le supportano (scalabilità, tolleranza ai guasti\ldots), con esempi pratici e influenza sul design e sulle prestazioni.

\textbf{Dove trovare la risposta}: sezione \emph{Caratteristiche dei Sistemi Distribuiti}, con la tabella caratteristica $\to$ tecnica.

\subsection*{Seconda sessione 2024 -- Informatica \#2}
\textbf{Traccia}: metodologie e modelli architetturali per applicazioni distribuite enterprise (microservizi, SOA, N-tier); impatto su scalabilità, tolleranza ai guasti, affidabilità; latenza, throughput, uso delle risorse; ruolo di cloud e virtualizzazione, con vantaggi, sfide e best practice.

\textbf{Dove trovare la risposta}: sezioni \emph{Architetture di Applicazioni Distribuite} e \emph{Cloud Computing e Virtualizzazione}.

\subsection*{Seconda sessione 2024 -- Informatica \#3}
\textbf{Traccia}: database per Big Data; modelli relazionale e non relazionale; scalabilità, flessibilità, velocità; elaborazione batch e streaming; impatto del teorema CAP; esempi e casi d'uso.

\textbf{Dove trovare la risposta}: sezioni \emph{Database SQL e NoSQL}, \emph{Big Data} e \emph{Il Teorema CAP}.

\subsection*{Prima sessione 2026 -- A2-INF-4}
\textbf{Traccia}: concetto di protocollo nei sistemi distribuiti, ruolo nell'interoperabilità, affidabilità e sincronizzazione; funzioni di un protocollo di comunicazione; confronto tra modelli.

\textbf{Dove trovare la risposta}: sezione \emph{Meccanismi Fondamentali}, e capitolo \emph{Reti di Calcolatori} per ISO/OSI e TCP/IP.

\subsection*{Prima sessione 2026 -- A3-INF-4 (tema di progetto)}
\textbf{Traccia}: sistema per un cinema multisala con interfaccia Web, backend applicativo, base di dati e modulo di Machine Learning; progetto del backend, delle API (metodi HTTP, endpoint, dati scambiati) e del flusso di prenotazione in UML.

\textbf{Dove trovare la risposta}: sezioni \emph{Architetture di Applicazioni Distribuite} e \emph{Meccanismi Fondamentali} (REST); capitoli \emph{Basi di Dati} e \emph{Progettazione del Software} per ER e UML.
